\chapter{Conclusiones generales}
\label{cap:conclusiones}

El objetivo general fue «Evaluar el desempeño de generalización de máquinas de aprendizaje profundo
frente al fenómeno de cambio de dominio en sistemas BCI multiusuario de
imaginación motora, mediante un protocolo de validación inter-sujeto que
permita la identificación de la arquitectura con mayor robustez». La evaluación permitió
caracterizar la degradación entre participantes y entre protocolos de adquisición, pero no identificar
una arquitectura de mayor robustez en todos los criterios. La elección dependió de la métrica,
el protocolo y las condiciones de evaluación.

\begin{enumerate}
  \item \textbf{Selección de datos (objetivo~\ref{oe:1}).} Se buscó un conjunto multisujeto
  de electroencefalografía superficial (sEEG) para imaginación motora. Se seleccionó
  \openbmi{} como conjunto principal y WBCIC-SHU para validación externa, con las clases
  de mano izquierda y derecha y los canales comunes. La evaluación incluyó $54$ participantes
  de \openbmi{} y $51$ de WBCIC-SHU; en este último se utilizó la primera sesión y los $45$
  canales comunes (sección~\ref{cap:resultados} y capítulo~\ref{cap:metodologia-datos}).

  \item \textbf{Preprocesamiento (objetivo~\ref{oe:2}).} Se requirió eliminar ruido y
  artefactos, tratar los canales y las bandas y preparar entradas compatibles. Se aplicaron
  filtrado, limpieza, segmentación y normalización. La selección de canales quedó implícita
  en las capas espaciales de cada arquitectura; se justificó no efectuar una selección explícita
  para evitar una decisión previa ajustada a los datos (capítulo~\ref{cap:metodologia-datos}).
  Las bandas se evaluaron como ablación: en el protocolo intra-sujeto, sobre $54$ participantes
  y con la semilla $2026$, la banda $\mu+\beta$ se asoció con $+0.0781$ de exactitud en el
  Conformer frente a su base. El contraste está condicionado al régimen de entrenamiento y
  las subbandas no se evaluaron en LaBraM ni EEGMamba
  (sección~\ref{sec:ablaciones-f4}).

  \item \textbf{Arquitecturas (objetivo~\ref{oe:3}).} Se implementaron y configuraron
  arquitecturas representativas para decodificar imaginación motora: Conformer y EEGMamba se
  entrenaron desde cero, mientras LaBraM Base recibió ajuste fino completo; \eegnet{} se
  incorporó como referencia. En validación dejando un sujeto fuera (\loso{}), sobre $54$
  participantes de \openbmi{} y con la media de las semillas $2026$, $2027$ y $2028$, las
  exactitudes fueron $0.747$, $0.715$, $0.755$ y $0.761$, respectivamente
  (sección~\ref{sec:comparacion}). Las semillas no compartieron completamente el régimen de
  optimización (sección~\ref{sec:semillas-por-arq}).

  \item \textbf{Evaluación comparativa (objetivo~\ref{oe:4}).} Se comparó la generalización
  inter-sujeto mediante exactitud y un indicador de degradación. El indicador de cambio de
  dominio propuesto (\dsi{}) comparó participantes vistos y no vistos con los mismos pesos,
  separando esa diferencia del volumen de entrenamiento. En \loso{}, con $54$ particiones
  de \openbmi{} y la media de tres semillas, fue de $+0.026$ a $+0.052$; solo LaBraM presentó
  degradación significativa tras Holm (sección~\ref{sec:dsi-resultados}). La regla aplicada,
  fijada después del primer análisis, prefirió el menor costo entre arquitecturas equivalentes
  en desempeño y degradación. Para exactitud \loso{}, sobre $54$ participantes y la media de
  tres semillas, la prueba bayesiana y TOST con Holm respaldaron la equivalencia por pares de
  \eegnet{}, Conformer y LaBraM dentro del margen de $\pm 0.03$. Para degradación, con las
  mismas condiciones, ambas pruebas respaldaron la equivalencia de los seis pares; entre
  \eegnet{} y LaBraM, la diferencia media fue $-0.021$ y el intervalo bootstrap del
  $90$\,\% fue $[-0.031, -0.012]$, próximo al límite del margen. Por su menor costo
  medido, \eegnet{} resultó preferida bajo esta regla. La decisión, aplicada al conjunto
  principal para las semillas y regímenes evaluados, no establece superioridad universal
  (secciones~\ref{sec:criterios-adicionales}, \ref{sec:comparacion} y
  \ref{sec:aplicacion-robustez}).
\end{enumerate}

La pregunta de investigación planteó cómo determinar la arquitectura de aprendizaje profundo
con mejor desempeño ante el cambio de dominio en interfaces cerebro-computador (BCI)
multiusuario de imaginación motora. La respuesta requiere comparar bajo particiones comunes,
medir la degradación y considerar la incertidumbre y el costo: en las condiciones estudiadas
se obtuvo una preferencia condicionada por \eegnet{}, sin evidencia para un ganador universal.
En la transferencia de modelos entrenados con $44$ participantes de \openbmi{} a WBCIC-SHU,
con la media de tres semillas, la exactitud cayó entre $0.095$ y $0.123$ respecto de los
$10$ participantes apartados del origen; la diferencia entre esa caída y la atribuida al
cambio de persona no fue significativa (sección~\ref{sec:descomposicion}).

La contribución principal consistió en distinguir la evaluación entre personas del cambio de
protocolo de adquisición y en proponer el \dsi{} para comparar sujetos vistos y no vistos
sin cambiar los pesos. El repositorio documenta el código y los procedimientos de evaluación
\cite{ayala2026repo}. La comparación común permitió reconocer ventajas relativas: la
equivalencia práctica en exactitud y degradación y el menor costo favorecieron a \eegnet{}
bajo la regla aplicada, fijada después del primer análisis para el conjunto principal y los
regímenes evaluados, mientras otros criterios de robustez favorecieron a otras arquitecturas
(sección~\ref{sec:aplicacion-robustez}). Las dificultades de interpretación incluyeron la
selección temprana del modelo intra-sujeto, los cambios de régimen entre semillas y la
precisión limitada de una sola partición de sujetos apartados. Las ablaciones se limitaron
a una semilla y el cambio de protocolo agrupó varios factores de adquisición, por lo que sus
causas individuales no quedaron establecidas (sección~\ref{sec:limitaciones}).

Respecto del anteproyecto, la comparación se concretó con \openbmi{} como conjunto principal,
WBCIC-SHU como validación externa y \eegnet{} como línea base adicional. El Conformer se
evaluó sin adaptación de dominio como componente principal; Conformer y EEGMamba se
entrenaron desde cero y LaBraM Base mediante ajuste fino completo. Se añadieron alineamiento
euclidiano y calibración con pocos ensayos, y la comparación estadística empleó Wilcoxon con
corrección de Holm, TOST y remuestreo bootstrap en lugar de la prueba de Friedman con
comparaciones posteriores prevista inicialmente. Estas decisiones concretaron la ejecución
de los objetivos.
